% Options for packages loaded elsewhere
\PassOptionsToPackage{unicode}{hyperref}
\PassOptionsToPackage{hyphens}{url}
%
\documentclass[
]{article}
\usepackage{amsmath,amssymb}
\usepackage{iftex}
\ifPDFTeX
  \usepackage[T1]{fontenc}
  \usepackage[utf8]{inputenc}
  \usepackage{textcomp} % provide euro and other symbols
\else % if luatex or xetex
  \usepackage{unicode-math} % this also loads fontspec
  \defaultfontfeatures{Scale=MatchLowercase}
  \defaultfontfeatures[\rmfamily]{Ligatures=TeX,Scale=1}
\fi
\usepackage{lmodern}
\ifPDFTeX\else
  % xetex/luatex font selection
\fi
% Use upquote if available, for straight quotes in verbatim environments
\IfFileExists{upquote.sty}{\usepackage{upquote}}{}
\IfFileExists{microtype.sty}{% use microtype if available
  \usepackage[]{microtype}
  \UseMicrotypeSet[protrusion]{basicmath} % disable protrusion for tt fonts
}{}
\makeatletter
\@ifundefined{KOMAClassName}{% if non-KOMA class
  \IfFileExists{parskip.sty}{%
    \usepackage{parskip}
  }{% else
    \setlength{\parindent}{0pt}
    \setlength{\parskip}{6pt plus 2pt minus 1pt}}
}{% if KOMA class
  \KOMAoptions{parskip=half}}
\makeatother
\usepackage{xcolor}
\setlength{\emergencystretch}{3em} % prevent overfull lines
\providecommand{\tightlist}{%
  \setlength{\itemsep}{0pt}\setlength{\parskip}{0pt}}
\setcounter{secnumdepth}{-\maxdimen} % remove section numbering
\ifLuaTeX
  \usepackage{selnolig}  % disable illegal ligatures
\fi
\IfFileExists{bookmark.sty}{\usepackage{bookmark}}{\usepackage{hyperref}}
\IfFileExists{xurl.sty}{\usepackage{xurl}}{} % add URL line breaks if available
\urlstyle{same}
\hypersetup{
  hidelinks,
  pdfcreator={LaTeX via pandoc}}

\author{}
\date{}

\begin{document}

\textbf{Anonimne mreže\\
}Jovan Tričković 150/2021\textbf{\hfill\break
}

\hypertarget{sadrux17eaj}{%
\subsection{Sadržaj}\label{sadrux17eaj}}

\protect\hyperlink{_Toc232183512}{\textbf{1. Uvod}
\protect\hyperlink{_Toc232183512}{3}}

\protect\hyperlink{kako-funkcioniux161u-vpn-i-tor}{\textbf{2. Kako
funkcionišu VPN i Tor}
\protect\hyperlink{kako-funkcioniux161u-vpn-i-tor}{4}}

\protect\hyperlink{_Toc232183514}{\textbf{2.1 VPN}
\protect\hyperlink{_Toc232183514}{4}}

\protect\hyperlink{_Toc232183515}{\textbf{2.2 Tor mreža}
\protect\hyperlink{_Toc232183515}{4}}

\protect\hyperlink{_Toc232183516}{\textbf{2.3 Onion adrese}
\protect\hyperlink{_Toc232183516}{5}}

\protect\hyperlink{tehnike-privatnosti-kod-tor-mreux17ee}{\textbf{3.
Tehnike privatnosti kod tor mreže}
\protect\hyperlink{tehnike-privatnosti-kod-tor-mreux17ee}{5}}

\protect\hyperlink{_Toc232183518}{\textbf{3.1 Operation security OpSec}
\protect\hyperlink{_Toc232183518}{6}}

\protect\hyperlink{_Toc232183519}{\textbf{3.2 Kako se tor brani od
fingerprinting a} \protect\hyperlink{_Toc232183519}{6}}

\protect\hyperlink{_Toc232183520}{\textbf{3.3 Deep packet inspection i
metode cenzure} \protect\hyperlink{_Toc232183520}{7}}

\protect\hyperlink{_Toc232183521}{\textbf{3.4 Kombinovanje tor mreže i
VPN usluga} \protect\hyperlink{_Toc232183521}{7}}

\protect\hyperlink{_Toc232183522}{\textbf{3.5 Decentralizovani VPN i
alternativne mreže} \protect\hyperlink{_Toc232183522}{8}}

\protect\hyperlink{pravni-i-etiux10dki-aspekti-tor-mreux17ee-i-vpn-usluga}{\textbf{4.
Pravni i etički aspekti tor mreže i VPN usluga}
\protect\hyperlink{pravni-i-etiux10dki-aspekti-tor-mreux17ee-i-vpn-usluga}{8}}\\
\textbf{5. Zaključak} 10

\textbf{\hfill\break
}

\protect\hypertarget{_Toc232183512}{}{}\textbf{1. Uvod}

U dobu veštačke inteligencije najbitniji resurs su podaci i samim tim
kompanije ulažu dosta novca u razne metode prikupljanja podataka.
Kompanije kao što su Facebook, Google, TikTok, a neke druge, poput
GoodRx, su delile lične podatke pacijenata već pomenutim kompanijama
{[}15{]}. Dodatni rezultat ovoga je i pojava samocenzure, gde se
korisnici suzdržavaju da istražuju određene informacije ili da ispolje
određena mišljenja jer osećaju da su praćeni. Ovo je još više slučaj u
državama kao što je Kina gde je politički pritisak mnogo veći. Uz pomoć
alata veštačke inteligencije, još lakše je prikupljati i klasifikovati
ove podatke, a neki alati mogu i da identifikuju korisnike samo na
osnovu njihovog stila pisanja {[}1{]}. Sve ovo čini alate za privatnost
i anonimnost sve više i više atraktivnim. Bitno je reći da privatnost i
anonimnost nisu ista stvar. Anonimnost znači da se ti podaci ne mogu
povezati sa određenim identitetom, a privatnost uključuje stvari kao što
su enkripcija i nemogućnost pristupa tim podacima.

Za ove potrebe postoje razni alati koji svaki pokriva svoj domen. Zato
korisnik ipak treba prvo da razmisli od čega hoće da se zaštiti, jer
jedan alat ne može da pokrije sve potrebe. Pričaćemo najviše o alatima
koji pružaju zaštitu na nivou IP adresa, kao što su Tor mreža i VPN
usluge, ali ćemo ukratko spomenuti i druge vidove zaštite.

IP adrese nisu jedini način da se poveže ponašanje sa korisnikom.
Postoje razne tehnike ``uzimanja otisaka'', tj. na engleskom
fingerprinting, gde često korišćenjem JavaScript-a sajt može da zaključi
neke stvari o korisniku i ako se dovoljno unikatnih osobina može
izdvojiti, može se izvršiti identifikacija {[}2, 4{]}. Neki tipovi
fingerprinting-a su:

\begin{itemize}
\item
  \textbf{Canvas HTML5 fingerprinting:} Koristeći Canvas HTML5 API, u
  zavisnosti od operativnog sistema i grafičke kartice, tekst u
  pretraživaču može da ima mikro razlike i samim tim može da služi kao
  unikatni ID {[}2, 4{]}.
\item
  \textbf{Audio fingerprinting:} Pretraživači mogu da se razlikuju u
  audio steku (audio stack) {[}4{]}.
\item
  \textbf{Font fingerprinting:} Instalirani fontovi na OS korisnika mogu
  da smanje skup korisnika koji imaju ili nemaju neki određeni font
  instaliran {[}2, 11{]}.
\item
  \textbf{Screen fingerprinting:} Veličina ekrana i dubina boja (color
  depth) mogu da se iskoriste za deanonimizaciju {[}2, 4{]}.
\item
  \textbf{Vreme i sistemski lokalitet (system locale):} Korišćenjem JS-a
  sajt može da dobije vremensku zonu i mesto korisnika {[}2, 4, 11{]}.
\end{itemize}

Svaki pretraživač takođe koristi kolačiće, tj. cookies zarad neophodne
funkcionalnosti sajta. U njima se sadrže npr. informacije o trenutnoj
login sesiji koje su bitne da ne bismo morali da se stalno logujemo
svaki put kada promenimo stranicu na istom sajtu i to su first-party
kolačići, ali takođe postoje i third-party kolačići koje sajtovi dele sa
kupcima i drugim sajtovima koji služe da bi pratili korisnika svuda na
internetu {[}3{]}.

Većina ljudi ne zna kako politike privatnosti funkcionišu i kako
kompanije koriste njihove lične podatke. Anketa od ``Pew Research
Center'' ukazuje i na to. 56\% uvek pritisne ``Slažem se'', a 67\% je
izjavilo da razumeju veoma malo o tome kako kompanije koriste njihove
podatke {[}14, 17{]}.

\hypertarget{kako-funkcioniux161u-vpn-i-tor}{%
\subsection{2. Kako funkcionišu VPN i
Tor}\label{kako-funkcioniux161u-vpn-i-tor}}

\protect\hypertarget{_Toc232183514}{}{}\textbf{2.1 VPN}

Jedan od načina da se korisnik zaštiti na nivou IP adresa je VPN.
Korisnik uspostavlja enkriptovanu point-to-point konekciju sa serverom.
Praktično, zamenjuje svoju IP adresu koju mu je dodelio ISP (Internet
Service Provider) sa adresom koju mu je dodelio VPN server. ISP samo
vidi da je korisnik uspostavio vezu sa VPN-om, ali ne zna ništa van
toga. Ovo može da bude dobro ako korisnik želi da se zaštiti od uticaja
ISP-a, ali je problem što se ovako ipak odgovornost samo prebacuje.
Korisnik i dalje mora da veruje centralizovanom entitetu, tj. VPN
provajderu. VPN provajder može da slobodno čita protok podataka. Tu sada
dolazi u obzir veoma bitan koncept koji se zove \emph{no-logs policy}
gde provajder tvrdi da ne loguje ni početne IP adrese ni protok koji se
dešava između izvora i odredišta. Ovo nije pravni termin i opet se
vraćamo na to da je dobar izbor provajdera i vera u njega bitna. Ipak
postoje 3rd-party kompanije koje vrše revizije (audits) i javno
objavljuju svoje rezultate, što služi kao ``dokaz'' da je provajder
stvarno ispunio svoje obećanje. Takođe, u zavisnosti od toga u kojoj je
državi registrovana, kompanija može da bude i pod pravnim pritiskom da
prikuplja podatke. Zato je bitno da VPN provajderi budu registrovani u
državama sa jačim zakonima o privatnosti.

\textbf{DNS leakage (Curenje DNS-a)} Kada korisnik pretražuje internet,
on ne koristi IP adrese, već imena domena. Zbog toga svaki domen mora da
se razreši u IP adresu kojoj pripada koristeći DNS server, obično od
ISP-a ili neke druge popularne opcije kao što je Google, itd. Tokom
korišćenja VPN-a, sav izlazni saobraćaj ide direktno kroz VPN server,
što uključuje i DNS zahteve koji bivaju razrešeni od strane VPN
provajdera. DNS leakage (curenje) se dešava kada zahtev, umesto da ide
kroz razrešivač od VPN provajdera, ide kroz podrazumevani razrešivač
koji je podešen na sistemu. Ovo znači da je moguće videti koji su domeni
posećeni, što ubija poentu uopšte korišćenja VPN-a.

\protect\hypertarget{_Toc232183515}{}{}\textbf{2.2 Tor mreža}

Tor je decentralizovana i enkriptovana mreža. S obzirom da je
decentralizovana, faktor poverenja je mnogo manji nego kod VPN-ova. Tor
mreža sadrži čvorove (nodes) koje pružaju volonteri. Klijent uspostavlja
3 node-a između sebe i destinacije, gde svaki node zna prethodnika i
sledbenika. Ovo znači da destinacija zna samo krajnji node, a krajnji
node zna samo za srednji node i destinaciju, pa ne može da zna ko je
poslao zahtev. Ovo se uspostavlja korišćenjem PFS (Perfect Forward
Secrecy), tj. klijent ne uspostavlja ceo put sa samo jednim
kriptografskim korakom. Prvo uspostavlja kratkoročni ključ sesije sa
početnim čvorom (entry guard) koristeći ``Diffie-Hellman rukovanje''.
Nakon što je ta veza uspostavljena, onda proširuje zahtev na middle node
i na kraju šalje zahtev ka exit node-u. Pošto svaki sloj ima svoj ključ
koji nestaje nakon što se konekcija ugasi, ovo čini da u slučaju da neki
od node-ova bude kompromitovan kasnije, nije moguće dekriptovati
prethodni internet saobraćaj. Po ovom modelu slojevanja je Tor i dobio
svoje ime ``The Onion Router'', gde svaki sloj štiti prethodni, kao kod
luka {[}12{]}. Ipak, Tor mreža je sporija od VPN-a i sajtovi cešće
blokiraju ip adrese čvorova nego VPN ip adrese, pa to treba isto uzeti u
obzir.

\textbf{Internet saobraćaj kroz Tor mrežu} Generalno, moguće je
analizirati dužinu paketa koji se šalju i time otkriti informacije o
prirodi tih paketa. Da bi se ovo izbeglo, Tor segmentiše saobraćaj u
ćelije, gde svaka ima 512 bajtova. Ovo čini bilo kakvo praćenje
(tracking) težim i jedan je od vidova odbrane koje Tor implementira
{[}12{]}. Tor je takođe originalno zamišljen da koristi leaky-pipe
tehnologiju koja bi omogućavala da middle node funkcioniše kao exit node
i time procesuira više multipleksovanih TCP protoka, ali se to u praksi
trenutno ne koristi {[}12{]}. Problem kod Tor mreže je to šta ako neko
ima pristup početnom i krajnjem node-u? U tom slučaju je moguće pratiti
saobraćaj kroz entry guard i kroz exit node. Pošto entry guard zna
klijentovu IP adresu, i onaj koji posmatra entry guard je isto zna. Uz
to zna odredišnu IP adresu. Korišćenjem ML modela i statistike, moguće
je uz pomoć metapodataka, saobraćaja, dužina prenosa itd. da se koreliše
korisnik sa tim saobraćajem. Ovo se zove \emph{end-to-end traffic
confirmation attack} i nije ga moguće potpuno zaobići, ali ideja je da
se smanji rizik tako što se smanjuje skup entry guard-ova {[}8, 12{]}.
Middle i exit nodovi se mnogo češće smenjuju, ali to nije slučaj sa
entry node-om. Pošto svako može da hostuje Tor node, to znači da
maliciozni učesnik može da bude deo mreže. Smanjenjem početnog skupa se
ta šansa smanjuje, jer je onda i manja šansa da naiđemo na maliciozni
node. Sledeći problem kod Tor mreže je exit node trust. Svaki node vrši
enkripciju, ali krajnji node mora da dekriptuje poruku kako bi je
isporučio odredišnoj IP adresi. Opet se vraćamo na problem volontera.
Veliki dobitak koji Tor mreža ima u poređenju sa VPN-om je to što se ne
mora verovati centralizovanom entitetu, ali je to u isto vreme i
nedostatak. Ako saobraćaj između klijenta i odredišta ne koristi
enkriptovani protokol kao HTTPS, i umesto toga koristi HTTP, i ako je
node maliciozan, moguće je logovati sav saobraćaj. Takođe je moguće
izvršiti tehniku ``SSL stripping'', gde napadač degradira
(downgrade-uje) HTTPS konekciju na HTTP konekciju i time može da čita
saobraćaj {[}8{]}.

\protect\hypertarget{_Toc232183516}{}{}\textbf{2.3 Onion adrese}

Internet bi ugrubo mogao da se podeli na 2 dela: na clear web, što
podrazumeva sve javne IP adrese kojima se normalno može pristupiti bez
korišćenja Tor mreže, i deep web koji koristi onion adrese i kojima se
može pristupiti samo kroz Tor mrežu. Krajnja odredišta kojima se
pristupa imaju onion adrese koje se završavaju sa ``.onion'' i zovu se
onion serveri. Ideja je sledeća:

\begin{enumerate}
\def\labelenumi{\arabic{enumi}.}
\item
  Onion server uspostavlja kola, tj. putanje ka više node-ova, i nazove
  ih introduction points. Onda uzima njihove javne ključeve i stavlja ih
  u decentralizovanu heš tabelu (DHT) {[}12{]}.
\item
  Kada klijent želi da se konektuje na onion server, on čita DHT i uzima
  jedan od introduction points kome šalje enkriptovan zahtev i takođe
  uzima neki drugi node za sebe kao rendezvous point {[}12{]}.
\item
  Server je obavešten kroz introduction point da klijent želi da
  komunicira sa njim. Server gradi put ka rendezvous point-u gde će se
  vršiti komunikacija {[}12{]}.
\end{enumerate}

Motivacija za ovo je sledeća:

\begin{itemize}
\item
  Nije neophodno koristiti DNS razrešivač, zato što onion adrese u svom
  imenu imaju javni ključ i zato onion adrese imaju glomazna imena
  {[}12{]}.
\item
  Korišćenjem 2 Tor node-a kao centra komunikacije, ni server ni klijent
  ne znaju za tuđu IP adresu {[}12{]}.
\end{itemize}

\hypertarget{tehnike-privatnosti-kod-tor-mreux17ee}{%
\subsection{3. Tehnike privatnosti kod tor
mreže}\label{tehnike-privatnosti-kod-tor-mreux17ee}}

\protect\hypertarget{_Toc232183518}{}{}\textbf{3.1 Operation security
OpSec}

Deanonimizacija se obično ne dešava zbog napadača, već zbog greške u
rukovanju sigurnosnim alatima. OpSec se odnosi na metode koje korisnik
treba da koristi da ne bi kompromitovao svoju anonimnost i privatnost
{[}8{]}. Česta greška bi bila da korisnik koristi lične podatke dok
koristi Tor mrežu ili da se uloguje na svoje lične naloge koje koristi
na clear web-u. Idealno bi bilo ne mešati naloge koji se koriste na
clear web-u sa onima koji se koriste na deep web-u. Drugi način
kompromitacije je pomoću dokumenata i fajlova. Svi fajlovi sadrže
metapodatke, od kojih neki mogu da budu lokacija, vreme kreiranja itd.
Takođe se može analizirati generalno ponašanje na stranici. Način
pomeranja miša, brzina i dužina kucanja, posećene stranice, vreme
provedeno na stranici itd., sve se to može iskoristiti da se korisnik
fingerprint-uje.

\protect\hypertarget{_Toc232183519}{}{}\textbf{3.2 Kako se tor brani od
fingerprinting a}

Sakrivanje IP adrese pomoću Tor mreže gubi smisao ako sajtovi i dalje
mogu da identifikuju specifičnu mašinu kroz napredne tehnike "uzimanja
otisaka" (browser fingerprinting). Standardni pretraživači (poput
Chrome-a ili Edge-a) šalju ogroman broj jedinstvenih parametara hardvera
i softvera kako bi optimizovali prikazivanje stranica. Pošto je Tor
pretraživač modifikovana verzija Firefox-a, to znači da ima visok nivo
opcija koje korisnik može da konfirugiše, što uključuje i plugine. Ipak,
Tor radi suprotno; on pokušava da učini da svi korisnici Tor-a izgledaju
identično. Ovaj koncept se naziva "efekat krda" (the herd effect)
{[}11{]}.

Neki od tehničkih mehanizama odbrane unutar Tor pretraživača su
{[}11{]}:

\begin{itemize}
\item
  \textbf{Letterboxing:} Kada korisnik promeni veličinu prozora
  pretraživača, rezolucija ekrana postaje unikatna. Da bi to sprečilo,
  Tor pretraživač koristi tehniku zvanu letterboxing. On dodaje sive
  ivice (margine) oko veb-stranice, prisiljavajući prozor u kome se
  renderuje sadržaj da se uvek drži standardizovanih dimenzija (npr.
  umnožaka od 200x100 piksela). Tada sajt vidi samo generičku rezoluciju
  koju dele milioni drugih Tor korisnika. Ipak, iako su dimenzije
  standardizovane, preporučuje se da se koristi default dimenzija ekrana
  i još više je bitno da se izbegava full screen opcija, jer će onda
  rezolucija Tor pretraživača odgovarati rezoluciji monitora.
\item
  \textbf{Standardizacija korisničkog agenta (User-Agent):} Bez obzira
  na to da li korisnik pokreće Tor na Mac-u, Linux-u ili Windows
  kompjuteru, pretraživač aplikacijama javlja potpuno isti niz podataka
  (User-Agent). Sajtovima se uvek prijavljuje ista verzija Firefox-a i
  generički operativni sistem.
\item
  \textbf{Blokiranje Canvas i Audio API-ja:} Kada sajt pokuša da
  iskoristi HTML5 Canvas za iscrtavanje skrivenih grafičkih elemenata
  radi identifikacije hardvera, Tor pretraživač automatski blokira taj
  zahtev ili vraća prazne/lažne podatke. Slično važi i za audio
  podsistem -- Web Audio API je modifikovan tako da ne dozvoljava
  precizna merenja performansi zvučne karte.
\item
  \textbf{Restrikcija sistemskih fontova:} Tor pretraživač ignoriše
  fontove koji su lokalno instalirani na operativnom sistemu korisnika i
  umesto toga koristi sopstveni set standardnih fontova, čime neutrališe
  font fingerprinting.
\item
  \textbf{Vremenska zona i jezik:} Svi zahtevi se automatski prevode na
  en-US jezički lokal, a sistemska vremenska zona se unutar pretraživača
  prisilno postavlja na UTC (Coordinated Universal Time), bez obzira na
  geografsku lokaciju korisnika.
\end{itemize}

\protect\hypertarget{_Toc232183520}{}{}\textbf{3.3 Deep packet
inspection i metode cenzure}

Iako VPN i Tor pružaju snažnu enkripciju, njihovo korišćenje nije
nevidljivo. U odredjenim državama, internet provajderi (ISP) i državni
cenzori koriste tehniku poznatu kao dubinska analiza paketa (Deep Packet
Inspection - DPI). Standardna analiza paketa gleda samo "zaglavlje"
(header) kako bi saznala izvor i destinaciju saobraćaja, dok DPI dodatno
analizira samu frekvenciju i metapodatke sadržaja koji se prenosi. Na
ovaj način, softver može da prepozna "potpis" (signature) OpenVPN
protokola, WireGuard-a ili specifičan šablon Tor saobraćaja, iako je sam
sadržaj enkriptovan i nečitljiv.

Kada se ovakav saobraćaj detektuje, on može biti veštački usporen
(throttling) ili potpuno blokiran na nivou provajdera. Da bi se ovo
zaobišlo, razvijeni su dodatni alati za prikrivanje (obfuscation). Kod
VPN-a se koriste obfuscated serveri koji maskiraju VPN saobraćaj tako da
izgleda kao običan, nasumičan HTTPS saobraćaj. Tor mreža koristi sličan
koncept kroz pluggable transports (kao što je obfs4 protokol, poznatiji
kao bridge), koji menja strukturu paketa i otežava cenzorima da
automatski prepoznaju i blokiraju pristup mreži bez gašenja celokupnog
HTTPS saobraćaja u državi {[}10{]}.

\protect\hypertarget{_Toc232183521}{}{}\textbf{3.4 Kombinovanje tor
mreže i VPN usluga}

Često se postavlja pitanje da li korisnici treba da koriste VPN i Tor
istovremeno kako bi postigli maksimalnu sigurnost. U praksi, postoje dva
glavna načina za kombinovanje ovih alata: Tor over VPN (Tor preko VPN-a)
i VPN over Tor (VPN preko Tor-a) {[}8{]}.

Kod postavke Tor preko VPN-a, korisnik se prvo povezuje na VPN server, a
zatim preko te zaštićene konekcije pokreće Tor pretraživač. Ova metoda
je znatno popularnija i jednostavnija za postavljanje. Njena glavna
prednost je to što internet provajder ne vidi da korisnik koristi Tor
mrežu, već vidi samo standardni VPN saobraćaj. Takođe, VPN server
sprečava Tor entry guard čvor da vidi pravu IP adresu korisnika, što
umanjuje rizik u slučaju kompromitovanja početnog čvora. Međutim,
korisnik i dalje mora da veruje VPN provajderu da ne beleži njegov
početni mrežni protok.

S druge strane, VPN preko Tor-a podrazumeva da se saobraćaj prvo rutira
kroz Tor mrežu, a tek onda izlazi na VPN server pre nego što dođe do
konačne destinacije. Ovo je tehnički dosta kompleksnije za podešavanje.
Prednost ovog pristupa je to što VPN provajder ne zna pravu IP adresu
korisnika (vidi samo IP adresu Tor izlaznog čvora), a istovremeno se
zaobilazi problem zlonamernih Tor izlaznih čvorova jer je saobraćaj
zaštićen VPN enkripcijom. Ipak, ova metoda drastično usporava internet
konekciju i ostavlja korisnika ranjivim na globalne analize saobraćaja
{[}8{]}. Takodje, generalna preporuka je da ovo za veliki procenat
``threat'' modela nije potrebno i da pogrešna konfiguracija može doneti
više štete nego benefita.

\protect\hypertarget{_Toc232183522}{}{}\textbf{3.5 Decentralizovani VPN
i alternativne mreže}

Zbog nedostataka koje imaju centralizovane VPN usluge, razvijaju nove
tehnologije koje pokušavaju da reše taj problem kod VPN-ova i probleme
brzine kod Tor mreže, gde je jedna od tih inovacija koncept
decentralizovanog VPN-a (dVPN). Za razliku od komercijalnog VPN-a gde
jedna kompanija poseduje i kontroliše sve servere, dVPN mreže se
oslanjaju na globalnu mrežu volontera koji iznajmljuju svoj višak
mrežnog protoka. Sistem funkcioniše slično Tor-u po pitanju
distribucije, ali se često oslanja na blockchain tehnologiju i pametne
ugovore (smart contracts) kako bi automatizovao kriptografsko rutiranje
saobraćaja {[}6{]}. Zbog toga, korisnici ne moraju da veruju centralnom
autoritetu da neće beležiti njihove logove, jer arhitektura sistema to
tehnički onemogućava, a ne postoji ni jedan centralni server koji bi
državne vlasti mogle da ugase.

Pored Tor-a, postoje i druge deepnet mreže koje pristupaju privatnosti
iz drugog ugla, poput I2P (Invisible Internet Project). I2P je izgrađen
isključivo kao zatvoreni ekosistem za peer-to-peer (P2P) komunikaciju
unutar same mreže, tj. ne može se koristiti za noramalan surface web
internet. Pošto I2P ne koristi izlazne čvorove ka spoljnom internetu, ta
ga čini znatno otpornijim na cenzuru i napade nadgledanja izlaznog
saobraćaja {[}5{]}.

\hypertarget{pravni-i-etiux10dki-aspekti-tor-mreux17ee-i-vpn-usluga}{%
\subsection{4. Pravni i etički aspekti tor mreže i VPN
usluga}\label{pravni-i-etiux10dki-aspekti-tor-mreux17ee-i-vpn-usluga}}

Tehnologije poput VPN-a i Tor mreže su često centar savremenih debata:
gde treba povući granicu između prava pojedinca na privatnost i potrebe
države da sprovodi zakone i štiti javnu bezbednost? Ovde jer reč je pre
o etičkom i pravnom problemu, zbog čega su anonimne mreže predmet
kontinuiranih rasprava među zakonodavcima, bezbednosnim službama,
aktivistima za ljudska prava i stručnjacima za sajber-bezbednost i samim
tim nije lako dati odgvor {[}18{]}.

\textbf{Etička dilema: privatnost kao pravo ili anonimnost kao rizik}\\
Zagovornici alata za anonimnost polaze od pretpostavke da je privatnost
osnovno ljudsko pravo. Iz te perspektive, korišćenje Tor-a ili VPN-a,
tj. privatnost, se suštinski ne razlikuje od slobode govora. Argument je
da građani ne moraju da opravdavaju svoju potrebu za privatnošću;
naprotiv, teret opravdanja trebalo bi da bude na onima koji žele da
nadziru njihove aktivnosti.

Ovakav stav posebno dobija na značaju u državama koje primenjuju strogu
internet cenzuru ili nadzor komunikacija. Za ljude kao što su novinari,
političke aktivisti, `whistleblowers', istraživače itd. anonimne mreže
mogu predstavljati jedini bezbedan način za pristup informacijama ili
komunikaciju bez straha od represije. U takvim okolnostima, ove
tehnologije često se posmatraju kao instrument zaštite građanskih
sloboda i slobode izražavanja. Prema jednoj od analiza, korisnici Tor
mreže u politički represivnijim državama značajno češće koriste mrežu za
pristup standardnim internet sadržajima i zaobilaženje cenzure nego za
pristup skrivenim servisima povezanim sa kriminalnim aktivnostima {[}9,
18{]}.

Sa druge strane, kritičari ovih tehnologija ukazuju da anonimnost može
otežati identifikaciju pojedinaca koji učestvuju u ozbiljnim krivičnim
delima. Tržišta narkotika, distribucija malvera, prevare, krađa podataka
i razmena materijala seksualnog zlostavljanja dece samo su neki od
primera aktivnosti koje su tokom godina identifikovane na anonimnim
mrežama {[}18{]}. Iz perspektive vlasti, potpuna anonimnost može
predstavljati značajnu prepreku istragama i krivičnom gonjenju
odgovornih lica.

\textbf{Tehnologija dvostruke namene}\\
Tor se klasnično ponaša kao klasični primer tehnologije dvostruke
namene. Ovakve tehnologije nisu ni inherentno dobre ni inherentno loše,
jer njihov etički status zavisi od načina na koji se koriste. Ista mreža
koja omogućava novinaru da anonimno komunicira sa izvorom može omogućiti
i kriminalcu da prikrije svoj identitet. Zbog toga se etička procena ne
može zasnivati isključivo na postojanju tehnologije, već i na njenim
dominantnim obrascima upotrebe {[}18{]}.

Istraživanja donekle osporavaju popularnu predstavu da se Tor koristi
prvenstveno za kriminalne aktivnosti. Studija objavljena u časopisu
Proceedings of the National Academy of Sciences procenjuje da približno
6,7\% korisnika Tor mreže na prosečan dan pristupa onion servisima koji
su češće povezani sa nelegalnim sadržajem {[}9{]}. Drugim rečima, velika
većina korisnika Tor-a koristi mrežu za pristup običnom internetu (clear
web), a ne za aktivnosti koje se tradicionalno povezuju sa „dark
web``-om. Autori takođe pokazuju da je procenat potencijalno
problematične upotrebe viši u politički slobodnim državama (oko 7,8\%)
nego u represivnijim režimima (oko 4,8\%), što sugeriše da u državama sa
visokim nivoom cenzure Tor često služi prvenstveno za zaštitu osnovnih
građanskih prava {[}9{]}.

\textbf{Pravni status anonimnih mreža}\\
Pravni položaj VPN i Tor tehnologija značajno se razlikuje od države do
države. U većini demokratskih zemalja korišćenje VPN usluga i Tor
pretraživača samo po sebi nije protivzakonito. U tim državama,
privatnost se smatra legitimnim interesom građana, a upotreba tehničkih
sredstava za njenu zaštitu uglavnom je dozvoljena što se razlikuje od
država koje sprovode intenzivnu kontrolu interneta. Watson u svojoj
komparativnoj pravnoj analizi pokazuje da se odnos prema anonimnim
mrežama značajno razlikuje između država poput Sjedinjenih Američkih
Država i država koje imaju strože mehanizme kontrole interneta,
uključujući Kinu, Saudijsku Arabiju i Ujedinjene Arapske Emirate. U
takvim sistemima anonimne mreže se često posmatraju kao potencijalna
pretnja državnom nadzoru, pa se njihova upotreba može ograničavati
tehničkim i pravnim sredstvima {[}7{]}.

Ovakve razlike proizlaze iz različitih pravnih i političkih prioriteta.
Dok jedni sistemi naglašavaju slobodu izražavanja i zaštitu privatnosti,
drugi prioritet daju državnoj bezbednosti, kontroli informacija ili
očuvanju društvene stabilnosti. Posledica toga je da ista tehnologija
može biti posmatrana kao sredstvo zaštite ljudskih prava u jednoj
državi, a kao bezbednosni problem u drugoj.

\textbf{Debata o „zadnjim vratima`` u enkripciji}\\
Jedna od kontroverznih tema u ovoj oblasti jeste pitanje da li bi države
trebalo da imaju mogućnost pristupa enkriptovanim komunikacijama kroz
„zadnja vrata`` (backdoors). Zagovornici ove ideje tvrde da bi takav
pristup omogućio efikasnije istrage terorizma, organizovanog kriminala i
drugih teških krivičnih dela {[}16{]}. Prema ovom argumentu, potpuno
neprobojne komunikacije stvaraju „zone nedodirljivosti`` u kojima
kriminalci mogu da deluju bez straha od nadzora {[}16{]}.

Protivnici ovog pristupa smatraju da je takvo rešenje tehnički i etički
problematično {[}13{]}. Ukoliko postoji mehanizam kojim legitimni organi
mogu da zaobiđu zaštitu sistema, tada postoji i mogućnost da isti
mehanizam jednog dana bude zloupotrebljen od strane kriminalaca, stranih
obaveštajnih službi ili drugih neovlašćenih aktera {[}13{]}. Iz tog
razloga, jedan od argumenata je da ne postoji način da se naprave
„zadnja vrata`` koja bi bila dostupna samo dobronamernim akterima.
Slabljenje enkripcije radi lakšeg sprovođenja zakona moglo bi
istovremeno da ugrozi bezbednost miliona legitimnih korisnika interneta
{[}13, 16{]}.\\
\strut \\
\textbf{Debata of verifikaciji godina.}\\
Aktuelna je i debata o skorašnjem talasu verifikacije godina (age
verification) za društvene mreže u nekim zapadnim državama kao što su
Australija i UK {[}19, 20{]}, što je privuklo veliku pažnju i debatu.
Naime, vlasti i pristalice ovih zakona citiraju razne studije koje su
pokazale negativni uticaj društvenih mreža na razvoj dece {[}23{]} i
time apeluju da društvene mreže počnu da verifikuju godine svojih
korisnika. Time bi korisnici iz tih zemalja morali da pokažu da
prevazilaze starosnu granicu, i time bi manje dece bile izloženo
negativnom uticaju društvenih mreža. Kritičari ovog zastupanja kažu da
ovo predstavlja veliki problem privatnosti i da se time stvara baza
ličnih podataka koja je veoma atraktivna za hakovanja tj „data
breaches``. {[}22{]} Dodatni argument je to da ovi zakoni nemaju
dovoljno pozitivnog efekta jer su neefikasni {[}24{]}. Način na koji su
VPN usluge povezane sa ovim pitanjem je to što popularnost VPN usluga u
državama porasla nakon ovih zakona {[}21{]}, jer korsnici ne žele da
rizikuju svoju privatnost samo da bi koristili društvene mreže, pa se
postavlja pitanje: „Da li je to legitimno korišćenje ovih usluga?{}`` i
da li će to negativno uticati na pravni status VPN-ova. Trenutno,
nijedna država nije zabranila upotrebu VPN usluga zbog ovih zakona.\\
\strut \\
\textbf{5. Zaključak\\
\strut \\
}Alati poput VPN-a i Tor mreže dobijaju na popularnosti kako načini
profilisanja postaju sve agresivniji. Kako raste sofisticiranost metoda
praćena, poput fingeprintinga i kolačića, tako i načini da se te metode
zaobiđu postaju jače. Generalno, i decentralizovane i centralizovana
rešenja imaju svoje mane i prednosti, pa je zato bitno da korsnik
izabere alat koji mu najviše odgovara.

Ipak, tehnička rešenja su efikasna samo onoliko koliko korisnik zna
njima da upravlja, jer bez striktne operativne bezbednosti (OpSec), čak
i najnaprednija enkripcija gubi smisao. Na kraju, ove tehnologije ne
pružaju samo privatnost i anonimnost, već i bitne metode bezbednosti u
autoritarnim državama.

\textbf{Reference\\
\strut \\
}{[}1{]} arXiv preprint "De-Anonymization at Scale via Tournament-Style
Attribution" (2026) \url{https://arxiv.org/html/2601.12407v2}

{[}2{]} P. Laperdrix, W. Rudametkin, B. Baudry "Browser Fingerprinting:
A Survey" (2019) \url{https://arxiv.org/pdf/1905.01051}

{[}3{]} Federal Trade Commission (FTC) "How Websites and Apps Collect
and Use Your Information"
\url{https://consumer.ftc.gov/articles/how-websites-apps-collect-use-your-information}

{[}4{]} Electronic Frontier Foundation (EFF) "Cover Your Tracks: Learn
How Trackers Work" \url{https://coveryourtracks.eff.org/learn}

{[}5{]} I2P Project "Introduction - I2P (Invisible Internet Project)"
\url{https://i2p.net/en/docs/overview/intro/}

{[}6{]} arXiv preprint "VPN0: A Privacy-Preserving Decentralized Virtual
Private Network" (2019) \url{https://arxiv.org/abs/1910.00159}

{[}7{]} A. Watson "The Tor Network: A Global Inquiry into the Legal
Status of Anonymity Networks" (2012)
\url{https://journals.library.wustl.edu/globalstudies/article/id/681/}

{[}8{]} Journal of Cyber Security and Mobility "Security and Traffic
Analysis of Anonymity Networks" (2019)
\url{https://journals.riverpublishers.com/index.php/JCSANDM/article/view/5981?articlesBySameAuthorPage=2}

{[}9{]} E. Jardine, A. Lindner, G. Owenson "The potential harms of the
Tor anonymity network cluster disproportionately in free countries"
(2020) \url{https://pmc.ncbi.nlm.nih.gov/articles/PMC7749358/}

{[}10{]} The Tor Project "Tor Project Support: Using Bridges"
\url{https://support.torproject.org/tor-browser/circumvention/using-bridges/}

{[}11{]} The Tor Project "Tor Project Support: Fingerprinting
Protections"
\url{https://support.torproject.org/tor-browser/features/fingerprinting-protections/}

{[}12{]} R. Dingledine, N. Mathewson, P. Syverson "Tor: The
Second-Generation Onion Router" (2004)
\url{https://svn-archive.torproject.org/svn/projects/design-paper/tor-design.html}\\
\strut \\
\url{https://svn-archive.torproject.org/svn/projects/design-paper/tor-design.html}\strut \\
\url{https://arxiv.org/pdf/1905.01051}

{[}13{]} DIGITALEUROPE "Encryption: Finding the balance between privacy,
security and lawful data access" (2017)
\url{https://www.digitaleurope.org/resources/encryption-finding-the-balance-between-privacy-security-and-lawful-data-access/}

{[}14{]} Federal Trade Commission (FTC) "FTC Recommends Congress Require
Data Broker Industry to be More Transparent and Give Consumers Greater
Control Over Their Personal Information" (2014)
\url{https://www.ftc.gov/news-events/news/press-releases/2014/05/ftc-recommends-congress-require-data-broker-industry-be-more-transparent-give-consumers-greater}

{[}15{]} Federal Trade Commission (FTC) "FTC Enforcement Action to Bar
GoodRx from Sharing Consumers' Sensitive Health Info for Advertising"
(2023)
\url{https://www.ftc.gov/news-events/news/press-releases/2023/02/ftc-enforcement-action-bar-goodrx-sharing-consumers-sensitive-health-info-advertising}

{[}16{]} U.S. Department of Justice (DOJ) "International Statement:
End-to-End Encryption and Public Safety" (2020)
\url{https://www.justice.gov/archives/opa/pr/international-statement-end-end-encryption-and-public-safety}

{[}17{]} Pew Research Center "How Americans View Data Privacy" (2023)
\url{https://www.pewresearch.org/internet/2023/10/18/how-americans-view-data-privacy}

{[}18{]} M. Al-Saleh et al. "Shedding Light on the Dark Corners of the
Internet: A Survey of Tor Research" (2018)
\url{https://www.researchgate.net/publication/323627387_Shedding_Light_on_the_Dark_Corners_of_the_Internet_A_Survey_of_Tor_Research}

{[}19{]} UK Parliament "Online Safety Act 2023" (2023)
\url{https://www.legislation.gov.uk/ukpga/2023/50/contents}\\
\strut \\
{[}20{]} Parliament of Australia "Online Safety Amendment (Social Media
Minimum Age) Act" (2024)
\href{https://www.legislation.gov.au/C2024A00127/asmade/text}{Online
Safety Amendment (Social Media Minimum Age) Act 2024 - Federal Register
of Legislation}\\
\strut \\
{[}21{]} Cybernews "Age verification laws drive surge in bypass
discussions and VPN downloads" (2026)
\href{https://www.google.com/search?q=https://cybernews.com/news/age-verification-laws-drive-surge-vpn-downloads/\&authuser=1}{https://cybernews.com/news/age-verification-laws-drive-surge-vpn-downloads/}\\
\strut \\
{[}22{]} Malwarebytes "{[}updated{]} Age verification vendor Persona
left frontend exposed, researchers say" (2026)
\url{https://www.malwarebytes.com/blog/news/2026/02/age-verification-vendor-persona-left-frontend-exposed}\\
\strut \\
{[}23{]} U.S. Department of Health and Human Services (HHS) "Social
Media and Youth Mental Health: The U.S. Surgeon General's Advisory"
(2023)
\url{https://www.hhs.gov/sites/default/files/sg-youth-mental-health-social-media-advisory.pdf}\\
\strut \\
{[}24{]} B. Nansen, H. Sandberg, L. Bliss, S. Cohney "From Phreaking to
Sneaking: Children\textquotesingle s Circumvention of Social Media Age
Verification Systems" (2026) \url{https://arxiv.org/abs/2605.00368}

\end{document}
